\documentclass[11pt,a4paper]{article}
\usepackage[T1]{fontenc}
\usepackage[utf8]{inputenc}
\usepackage{amsmath,amssymb,amsthm}
\usepackage[margin=1in]{geometry}
\usepackage[colorlinks=true,linkcolor=blue,urlcolor=blue]{hyperref}
\theoremstyle{plain}
\newtheorem{theorem}{Theorem}
\newtheorem{lemma}[theorem]{Lemma}
\newtheorem{proposition}[theorem]{Proposition}
\newtheorem{corollary}[theorem]{Corollary}
\theoremstyle{definition}
\newtheorem{remark}[theorem]{Remark}

\title{The empirical recurrence for OEIS A200892, proved from an exact model}
\author{Adrian Perez Fontelles\\ \small Independent researcher}
\date{13 September 2026}
\begin{document}
\maketitle

\begin{abstract}
OEIS A200892 carries an empirical recurrence of order $10$. It is true, and it is
decidable rather than empirical. The entry's sequence is given exactly by an Ehrhart quasi-polynomial, from
which $a$ provably satisfies a MONIC linear recurrence of order $S=11$, derived below and not
assumed. The residual of the conjectured recurrence satisfies the same one, so whether it
annihilates $a$ is settled by a finite exact computation. No transfer matrix is involved and
the count is not a walk.
\end{abstract}

\noindent\small 2020 Mathematics Subject Classification. 05A15, 05A19, 11P21.\normalsize

\section{The conjecture}

OEIS A200892 is ``Number of 0..n arrays x(0..8) of 9 elements without any interior element greater than both neighbors.''. It has offset $1$ and begins
\[
200, 4059, 34350, 181336, 710976, 2269938, 6233356, 15250675,\ \dots
\]
The entry states, as a conjecture contributed by its author and never marked settled:
\begin{quote}
Empirical: a(n) = (2/2835)*n\^{}9 + (131/630)*n\^{}8 + (2803/945)*n\^{}7 + (1349/90)*n\^{}6 + (41449/1080)*n\^{}5 + (20423/360)*n\^{}4 + (1149293/22680)*n\^{}3 + (22741/840)*n\^{}2 + (2011/252)*n + 1.
\end{quote}
As of the ``Last modified'' line on the live entry (Oct 16 2017 12:23:26, revision 10)
this is still recorded as empirical, and nothing on the entry records it as proved.

\section{The count is a quasi-polynomial}

The array has a FIXED number $L$ of entries and an alphabet that grows with $n$. Every
condition the entry imposes is a boolean combination of statements
$\sum_i c_ix_i + c\,n \bowtie 0$ with integer coefficients and $\bowtie$ one of
$=,\ne,<,\le,>,\ge$: the bounds $|x_i|\le n$ or $0\le x_i\le n$, a vanishing sum, a difference
that must not vanish, a pair that must total exactly $n$. Each is homogeneous in $(x,n)$
jointly, so the admissible $x$ for a given $n$ are precisely the integer points at height $n$
of a finite union of relatively open rational cones in $\mathbb{R}^{L+1}$, and $a(n)$ is
the Ehrhart quasi-polynomial of that union: of degree at most $L$, with period dividing the
heights of the cones' ray generators.

\section{The bound}

The period is derived, not assumed. Every ray of every cell of the arrangement is cut out by
$L$ linearly independent hyperplanes taken from the constraint set; writing that $L\times(L+1)$
system as $M$, its null direction has $t$-component $\pm\det$ of the $x$-parts, so the
primitive generator $(x^*,t^*)$ has $t^*$ dividing that determinant. Let $T$ be the set of
those determinants. A simplicial cone in a triangulation has at most $L+1$ rays, each
contributing a factor $1-z^{t_i}$ with $t_i\in T$, so
\[
A(z)\;=\;\prod_{d\,\mid\,t\ \text{for some}\ t\in T}\Phi_d(z)^{L+1}
\]
annihilates $a$, and $S=\deg A=11$. A congruence condition modulo $m$ is counted one residue
class at a time; a class is a coset of $m\mathbb{Z}^L$, and a generator primitive in
$\mathbb{Z}^{L+1}$ must be scaled by a divisor of $m$ to lie in it, so each $t$ is replaced
by $mt$ and nothing else changes.

Over the common denominator the numerator has degree below that of $A$ for every cell of the
arrangement that HAS a ray, because such a cell's numerator collects a fundamental
parallelepiped whose heights are below the sum of its generators' heights. The origin is a cell
too and it has none: its series is the constant $1$, which over $A$ is $A/A$ and reaches degree
exactly $\deg A$. The bound in force is therefore $zA(z)$, of order $S=11$, and
$a(0),\dots,a(S-1)$, computed exactly, determine every later term. Nothing is fitted, and the
annihilator is tested on terms the model was not asked for before anything is claimed --- a
bound one too small fails that test, which is how this one was found.

\section{The proof}

\begin{lemma}\label{lem:crit}
Let $A(z)=z^S-\sum_{j<S}\alpha_jz^{j}$ be the monic annihilator derived above and
$q(z)=z^{10}-\sum_ic_iz^{10-i}$ the characteristic polynomial of the conjectured
recurrence. The residual $r(n)=a(n)-\sum_ic_ia(n-i)$ is a linear combination of shifts of $a$,
so $A$ annihilates $r$ as well. Since $A(0)\ne0$ the recurrence it gives runs in both
directions, and $S$ consecutive zeros of $r$ force $r$ to vanish at every later index.
\end{lemma}

\begin{proof}
$A$ annihilates $a$, hence every shift of $a$, hence any linear combination of shifts; that is
$r$. A monic recurrence with nonzero constant term determines each term from the $S$ before it
and each from the $S$ after it, so a run of $S$ zeros propagates.
\end{proof}

\begin{theorem}
\[
a(n)\;=\;10\,a(n-1) - 45\,a(n-2) + 120\,a(n-3) - 210\,a(n-4) + 252\,a(n-5) - 210\,a(n-6) + 120\,a(n-7) - 45\,a(n-8) + 10\,a(n-9) - a(n-10)
\]
for every $n>0$.
\end{theorem}

\begin{proof}
The terms $a(0),\dots$ were computed exactly in integer arithmetic from the model, extended by
$A$ where the model itself was not evaluated, and $r(n)$ formed for each index. Those integers
vanish from $n=0+1$ onwards, and the run exceeds $S+10$, which by
Lemma~\ref{lem:crit} settles every larger index at once.
\end{proof}

\section{Verification}

Three checks, each independent of the algebra above.

First, the model reproduces all $22$ terms the entry publishes, exactly, in integer
arithmetic. This is what ties the model to the entry: it is built by reading the entry's
English, and a misreading gives different counts.

Second, the conjectured recurrence was evaluated directly on those published terms, with no
model involved, and holds wherever the proved range and the data overlap.

Third, the derived annihilator $A$ was itself tested on terms it had not been given: the model
was evaluated beyond the $S$ values that determine it and $A$ reproduced them. A bound that is
too small fails this test, which is why it is run.

\begin{thebibliography}{9}
\bibitem{oeis} The OEIS Foundation, \emph{The On-Line Encyclopedia of Integer
Sequences}, \url{https://oeis.org}, 2026. Sequence A200892.
\bibitem{beckrobins} M.~Beck and S.~Robins, \emph{Computing the Continuous Discretely},
2nd ed., Springer, 2015. (Ehrhart quasi-polynomials and their periods.)
\bibitem{stanley} R.~P.~Stanley, \emph{Enumerative Combinatorics}, Volume 1, 2nd ed.,
Cambridge University Press, 2011.
\bibitem{ds} F.~Diamond and J.~Shurman, \emph{A First Course in Modular Forms},
Springer GTM 228, 2005.
\end{thebibliography}

\end{document}
